% ============================================================
% preambulo-contrato.tex
% ------------------------------------------------------------
% CONTRATO VERSIÓN 1
% ------------------------------------------------------------
% DEFINE TODO LO QUE docbook-to-latex.xsl PUEDE EMITIR EN LOS
% FRAGMENTOS DE CAPÍTULO. ES LA CAPA QUE UNA EDITORIAL
% CONSERVA CUANDO REEMPLAZA LA ESTÉTICA: SI FALTA UNA DE
% ESTAS DEFINICIONES, EL LIBRO NO COMPILA; SI ESTÁ DEFINIDA
% DISTINTO, COMPILA Y SALE MAL.
%
% CUANDO SE AGREGUE VOCABULARIO NUEVO, SUBIR LA VERSIÓN DEL
% CONTRATO. EL main.tex DECLARA LA QUE NECESITA Y ESTA CAPA
% LA VERIFICA: ASÍ UN PREÁMBULO VIEJO FRENA EN LA PRIMERA
% LÍNEA CON UN MENSAJE ÚTIL, EN VEZ DE DAR UN
% "Undefined control sequence" EN LA PÁGINA 140.
% ------------------------------------------------------------
% DEPENDE DE preambulo-proyecto.tex, QUE VA ANTES.
% ============================================================

\def\gbContratoVersion{1}

% ============================================================
% \gbRequiereContrato{N}
% PROPÓSITO : EL main.tex DECLARA QUÉ VERSIÓN NECESITA.
% PARÁMETRO : N — versión requerida
% ============================================================
\newcommand{\gbRequiereContrato}[1]{%
  \ifnum#1>\gbContratoVersion\relax
    \PackageError{gbpublisher}{%
      Este preambulo implementa el contrato \gbContratoVersion\space
      y el libro requiere el #1}{%
      Actualizar preambulo-contrato.tex.}%
  \fi}

% ============================================================
% HOJAS DE CORTESÍA
% ------------------------------------------------------------
% Dos páginas en blanco al inicio —una hoja— que todo libro
% lleva. Entran en la paginación: son i y ii, y el sumario
% empieza en iii.
%
% \hbox{} ES IMPRESCINDIBLE: UNA PÁGINA SIN NINGÚN MATERIAL NO
% SE EMITE, ASÍ QUE UN \newpage SOBRE UNA PÁGINA VACÍA NO
% PRODUCE NADA Y EL CONTADOR NO AVANZA. LA CAJA VACÍA DA EL
% MÍNIMO DE CONTENIDO QUE OBLIGA A COMPONERLA.
%
% \thispagestyle{empty} LES QUITA EL FOLIO: SE CUENTAN PERO NO
% SE IMPRIME EL NÚMERO.
% ============================================================
\newcommand{\gbHojasCortesia}{%
	\thispagestyle{empty}\hbox{}\newpage
	\thispagestyle{empty}\hbox{}\newpage}

% ============================================================
% ENCABEZADO DE PIEZA GENERADA
% ------------------------------------------------------------
% Para las piezas sin cuerpo que llevan título propio. El título
% sale de capitulos.titulo_capitulo, que es donde el editor lo
% escribe, y no del preámbulo.
%
% \markboth actualiza el folio corrido, que \chapter* no toca.
% ============================================================
\newcommand{\gbEncabezadoPieza}[1]{%
	\chapter*{#1}%
	\addcontentsline{toc}{chapter}{#1}%
	\markboth{\MakeUppercase{#1}}{\MakeUppercase{#1}}}

% ============================================================
% ÍNDICE CON SU ENCABEZADO
% ------------------------------------------------------------
% \printindex HACE UN \clearpage INTERNO AUNQUE EL ÍNDICE SE
% DECLARE CON title={}. SIN NEUTRALIZARLO, EL ENCABEZADO QUEDA
% SOLO EN UNA PÁGINA Y LA LISTA EMPIEZA EN LA SIGUIENTE.
%
% EL \let VA DENTRO DE UN GRUPO: FUERA DE ÉL \clearpage TIENE QUE
% SEGUIR FUNCIONANDO, Y EL COLOFÓN QUE VIENE DESPUÉS LO NECESITA.
% ============================================================
\newcommand{\gbIndice}[2]{%
	\gbEncabezadoPieza{#1}%
	\begingroup
	\let\clearpage\relax
	\raggedright
	\small
	\printindex[#2]%
	\endgroup}

% ============================================================
% 1. PAQUETES QUE EL CONTRATO NECESITA
% ------------------------------------------------------------
% SOLO LO QUE HACE FALTA PARA QUE UN FRAGMENTO COMPILE. LA
% ESTÉTICA CARGA LO SUYO APARTE. TODOS ESTÁN EN TeX Live Y
% NINGUNO PIDE --shell-escape.
% ============================================================
\usepackage{etoolbox}
\usepackage{graphicx}
\usepackage{booktabs}
\usepackage{array}
\usepackage{listings}
\usepackage{ragged2e}
\usepackage{xcolor}
\usepackage[most]{tcolorbox}

% ------------------------------------------------------------
% IDIOMA
% ------------------------------------------------------------
% babel VA EN EL CONTRATO Y NO EN LA ESTÉTICA, POR DOS RAZONES.
% LA PRIMERA ES DE ORDEN: biblatex PIDE CARGARSE DESPUÉS DE
% babel Y csquotes, Y LA ESTÉTICA VA DESPUÉS DE ESTA CAPA.
% LA SEGUNDA ES DE FONDO: EL IDIOMA DECIDE LAS COMILLAS QUE
% PRODUCE \enquote, LA PARTICIÓN DE PALABRAS Y LAS CADENAS DE
% LA BIBLIOGRAFÍA. SON GARANTÍAS DEL CONTRATO, NO DECISIONES
% DE DISEÑO: UNA EDITORIAL QUE REEMPLAZA LA ESTÉTICA NO DEBE
% PODER CAMBIAR EL IDIOMA DEL LIBRO.
\usepackage[\gbidiomas]{babel}
\frenchspacing

% COMILLAS SEGÚN EL IDIOMA ACTIVO: EN ESPAÑOL DA « ».
% EL XSLT EMITE \enquote Y NUNCA COMILLAS LITERALES.
\usepackage[autostyle=true]{csquotes}


% ============================================================
% 3. BIBLIOGRAFÍA
% ------------------------------------------------------------
% EL ARCHIVO biblatex-<estandar>-config.tex ES RESPONSABLE DE
% CARGAR biblatex CON EL ESTILO QUE CORRESPONDA. SON NUESTROS
% ARCHIVOS, VIAJAN EN EL PAQUETE Y SE COPIAN AL PROYECTO: NO
% SE DEPENDE DE NINGÚN ESTILO REGISTRADO EN TeX Live.
%
% LA RAMA bibtex (vancouver) NO ESTÁ IMPLEMENTADA. EL XSLT
% CORTA ANTES DE GENERAR SI EL ESTILO LA REQUIERE.
% ============================================================
\input{cita-\estandar-config.tex}

% ============================================================
% Índice de autores del aparato bibliográfico
% ------------------------------------------------------------
% Lo llena biblatex: con indexing=true en el estilo, cada nombre
% citado y cada nombre de la bibliografía se indexan solos. El
% editor no escribe ningún \index.
%
% esindex ORDENA LOS APELLIDOS COMPUESTOS CASTELLANOS COMO
% CORRESPONDE: «de la Torre» SE ALFABETIZA POR Torre PERO SE
% IMPRIME CON LA PARTÍCULA. SIN ÉL, EL ORDEN POR OMISIÓN LOS
% MANDA A LA D.
%
% VA EN EL CONTRATO PORQUE ES LO QUE HACE QUE LA PIEZA
% indice_autores_bibliografia DIGA LA VERDAD. EL \makeindex ES
% DISEÑO Y VA EN LA ESTÉTICA.
% ============================================================
\usepackage{esindex}
\DeclareIndexNameFormat{default}{%
	\usebibmacro{index:name}{\esindex[names]}
	{\namepartfamily}
	{\namepartgiven}
	{\namepartprefix}
	{\namepartsuffix}}

% ============================================================
% 4. PÁGINAS Y SALTOS
% ============================================================

% ------------------------------------------------------------
% \gbclearrecto — SALTA A LA PRÓXIMA PÁGINA IMPAR
% \gbclearverso — SALTA A LA PRÓXIMA PÁGINA PAR
% ------------------------------------------------------------
% LA PÁGINA DE RELLENO LLEVA \hbox{} PORQUE UNA PÁGINA SIN
% CONTENIDO NO SE EMITE. NO SE USA UN CARÁCTER BLANCO: ESO
% DEJA TEXTO EXTRAÍBLE EN EL PDF, QUE APARECE AL COPIAR, EN
% pdftotext Y EN UN LECTOR DE PANTALLA.
% EL CONTADOR SIGUE CORRIENDO: SOLO SE OCULTA EL FOLIO.
% ------------------------------------------------------------
\makeatletter
\newcommand{\gbclearrecto}{%
  \clearpage
  \ifodd\c@page\else
    \thispagestyle{empty}\hbox{}\newpage
  \fi}
\newcommand{\gbclearverso}{%
  \clearpage
  \ifodd\c@page
    \thispagestyle{empty}\hbox{}\newpage
  \fi}
% LA BLANCA QUE GENERA \cleardoublepage SALE, POR OMISIÓN, CON
% EL FOLIO Y EL ENCABEZADO CORRIENTES. ESTE PARCHE LA DEJA
% VACÍA, QUE ES LO QUE SE ESPERA DE UNA PÁGINA EN BLANCO.
\renewcommand{\cleardoublepage}{\gbclearrecto}
\makeatother

% ------------------------------------------------------------
% gbpaginaespecial / gbpaginaimagen
% ------------------------------------------------------------
% PIEZAS SUELTAS DE PRELIMINARES: UNA DEDICATORIA, UN
% RECORDATORIO, UNA FOTO A PÁGINA ENTERA. VAN EN RECTO, SIN
% FOLIO NI ENCABEZADO, Y SU VUELTA QUEDA EN BLANCO. NO LLEVAN
% \chapter Y NO ENTRAN AL ÍNDICE.
% LA DE IMAGEN ROMPE LA CAJA DE TEXTO; LA OTRA LA RESPETA.
% ------------------------------------------------------------
\newenvironment{gbpaginaespecial}
  {\gbclearrecto\thispagestyle{empty}\begin{center}}
  {\end{center}\gbclearrecto}

\newenvironment{gbpaginaimagen}
  {\gbclearrecto\thispagestyle{empty}%
   \begingroup\centering\setlength{\parindent}{0pt}}
  {\par\endgroup\gbclearrecto}

% ------------------------------------------------------------
% gbcolofon
% ------------------------------------------------------------
% VA EN VERSO. NO PUEDE SER UN \chapter*: EN LA CLASE book,
% \chapter FUERZA EL RECTO PORQUE @openright ESTÁ ACTIVO, Y
% CUALQUIER SALTO PREVIO QUEDA ANULADO.
% ------------------------------------------------------------
\newenvironment{gbcolofon}
  {\gbclearverso\thispagestyle{empty}\justifying}
  {\clearpage}

% ============================================================
% 5. ESTRUCTURA
% ============================================================

% ------------------------------------------------------------
% \gbContadoresSinCapitulo
% ------------------------------------------------------------
% SE LLAMA UNA VEZ AL ENTRAR EN PRELIMINARES Y UNA VEZ AL
% ENTRAR EN POSLIMINARES. NO POR PIEZA.
%
% EN ESAS DOS ZONAS \chapter NO EJECUTA \refstepcounter{chapter}
% (book.cls, \@chapter, RAMA \if@mainmatter). ESE COMANDO ES LO
% ÚNICO QUE INCREMENTA EL CONTADOR DE CAPÍTULO Y, A LA VEZ,
% REINICIA LOS QUE DEPENDEN DE ÉL: section, footnote, figure,
% table Y equation. SI NO CORRE, NO PASA NINGUNA DE LAS DOS.
%
% LLEVAR chapter A CERO ACTIVA LA GUARDA \ifnum \c@chapter>\z@
% DE \thefigure, \thetable Y \theequation, QUE OMITE EL PREFIJO.
% SIN ESTO, UNA FIGURA DE POSLIMINARES SALE "Figura 1.2", CON
% EL NÚMERO DEL ÚLTIMO CAPÍTULO DEL CUERPO. VERIFICADO.
%
% LAS FIGURAS Y TABLAS NO SE REINICIAN POR PIEZA A PROPÓSITO:
% EN PRELIMINARES LA NUMERACIÓN CORRIDA (1, 2, 3) ES LA CORRECTA
% Y REINICIAR PRODUCIRÍA DOS "Figura 1" CON \label EN COLISIÓN.
% LAS NOTAS AL PIE SÍ VAN POR PIEZA: \thefootnote ES ARÁBIGO
% PELADO Y NO HAY PREFIJO QUE COLISIONE.
% ------------------------------------------------------------
\newcommand{\gbContadoresSinCapitulo}{%
  \setcounter{chapter}{0}%
  \setcounter{section}{0}%
  \setcounter{figure}{0}%
  \setcounter{table}{0}%
  \setcounter{equation}{0}%
}

% ------------------------------------------------------------
% \gbPiezaSinNumerar
% ------------------------------------------------------------
% PARA UNA PIEZA SIN NUMERAR QUE VA EN LA ZONA DE CUERPO. EL
% CASO TÍPICO SON LAS CONCLUSIONES, QUE DEBEN PRECEDER A
% \appendix PARA QUE LOS APÉNDICES NUMEREN CON LETRA.
%
% COMO ESAS PIEZAS SE EMITEN CON \chapter*, NO CORRE
% \refstepcounter{chapter} Y NO SE REINICIA NADA: SU PRIMERA
% NOTA AL PIE ARRANCA DONDE LA DEJÓ EL CAPÍTULO ANTERIOR.
% VERIFICADO.
%
% NO SE TOCA EL CONTADOR DE CAPÍTULO: EL APÉNDICE POSTERIOR LO
% NECESITA INTACTO, Y \appendix YA LO LLEVA A CERO POR SU
% CUENTA. TAMPOCO SE TOCAN figure NI table: ESAS PIEZAS NO
% LLEVAN NI CUADROS NI FIGURAS.
% ------------------------------------------------------------
% TAMBIÉN APAGA LA NUMERACIÓN DE SECCIONES. SIN ESO, LA
% PRIMERA SECCIÓN DE LAS CONCLUSIONES SALE "1.1", QUE COLISIONA
% CON LA SECCIÓN 1.1 REAL DEL CAPÍTULO 1: \thesection ES
% \thechapter.\arabic{section} Y EL CONTADOR DE CAPÍTULO SIGUE
% EN 1. VERIFICADO EN EL SUMARIO.
% LA ZONA DE APÉNDICES LA VUELVE A ENCENDER.
\newcommand{\gbPiezaSinNumerar}{%
  \setcounter{footnote}{0}%
  \setcounter{section}{0}%
  \setcounter{secnumdepth}{0}%
}

% ------------------------------------------------------------
% \gbParte{título}{subtítulo}
% ------------------------------------------------------------
% LAS PARTES NO SON ARCHIVOS: SON METADATOS DE partes_libro Y
% LAS EMITE GAMBAS EN EL main.tex. NO LLEVAN TEXTO PROPIO.
% EL SUBTÍTULO PUEDE VENIR VACÍO. AL ÍNDICE VA SOLO EL TÍTULO.
% ------------------------------------------------------------
\newcommand{\gbParte}[2]{%
  \part[#1]{#1%
    \ifstrempty{#2}{}{\\[0.4em]{\normalsize\mdseries\itshape #2}}%
  }}

% ------------------------------------------------------------
% \gbCapitulo{título}{subtítulo}
% ------------------------------------------------------------
% EQUIVALENTE PARA CAPÍTULOS CON SUBTÍTULO.
% ------------------------------------------------------------
\newcommand{\gbCapitulo}[2]{%
  \chapter[#1]{#1%
    \ifstrempty{#2}{}{\\[0.3em]{\normalsize\mdseries\itshape #2}}%
  }}

% ============================================================
% 6. BLOQUES
% ============================================================

% CITA EN BLOQUE CON ATRIBUCIÓN OPCIONAL.
% NO ES UN EPÍGRAFE: DOCBOOK LOS DISTINGUE Y EL XSLT TAMBIÉN.
\newenvironment{gbcita}
  {\begin{quote}\small}
  {\end{quote}}
\newcommand{\gbatribucion}[1]{\par\vspace{2pt}\hfill\textup{\small#1}\par}

% EPÍGRAFE. LA ATRIBUCIÓN PUEDE SER UNA CITA BIBLIOGRÁFICA:
% EL XSLT RESUELVE UN <biblioref> DE ADENTRO A \textcite.
\newcommand{\gbepigrafe}[2]{%
  \par\vspace{0.5em}%
  \begin{flushright}%
  \begin{minipage}{0.72\linewidth}%
    \small\itshape #1\par
    \vspace{2pt}\raggedleft\upshape\footnotesize #2\par
  \end{minipage}%
  \end{flushright}%
  \vspace{0.8em}\par}

% RECUADRO LATERAL. UN ARGUMENTO: info | supplementary.
\newenvironment{gbsidebar}[1]
  {\begin{tcolorbox}[colback=black!4,colframe=black!4,
     boxrule=0pt,arc=2pt,left=6pt,right=6pt,top=4pt,bottom=4pt,
     before skip=6pt,after skip=6pt]\small\sffamily}
  {\end{tcolorbox}}

% VERSO. EL XSLT YA EMITE \\ ENTRE LÍNEAS.
\newenvironment{gbverso}
  {\begin{verse}\itshape}
  {\end{verse}}

% PARLAMENTO Y LOCUTOR. DOCBOOK 5.2 BASE NO TIENE dialogue NI
% drama: EL CANÓNICO USA para role="speech" CON
% emphasis role="speaker" (RC-DB-04).
\newenvironment{gbparlamento}
  {\par\begingroup\setlength{\parindent}{0pt}\hangindent=1.2em\hangafter=1}
  {\par\endgroup}
\newcommand{\gblocutor}[1]{\textsc{\MakeLowercase{#1}}\quad}

% ============================================================
% 7. LLAMADAS DE ATENCIÓN
% ============================================================
% CADA LLAMADA DE ATENCIÓN ES UN tcolorbox CON TÍTULO FIJO.
% NO SE USA UN .style DE pgfkeys PARAMETRIZADO: EL #1 DE UN
% \newcommand DENTRO DE UN .style LO CONSUME pgfkeys Y DA
% "You already have nine parameters", QUE NO DICE NADA SOBRE
% LA CAUSA. CON UN ENTORNO POR CLASE EL PROBLEMA NO EXISTE.
\tcbset{gbadmon/.style={
  colback=black!3, colframe=black!25,
  boxrule=0.4pt, arc=1pt,
  left=6pt, right=6pt, top=4pt, bottom=4pt,
  before skip=6pt, after skip=6pt,
  fonttitle=\bfseries\sffamily\small}}

\newenvironment{gbnota}
  {\begin{tcolorbox}[gbadmon,title={Nota}]\small}{\end{tcolorbox}}
\newenvironment{gbconsejo}
  {\begin{tcolorbox}[gbadmon,title={Sugerencia}]\small}{\end{tcolorbox}}
\newenvironment{gbadvertencia}
  {\begin{tcolorbox}[gbadmon,title={Advertencia}]\small}{\end{tcolorbox}}
\newenvironment{gbimportante}
  {\begin{tcolorbox}[gbadmon,title={Importante}]\small}{\end{tcolorbox}}
\newenvironment{gbprecaucion}
  {\begin{tcolorbox}[gbadmon,title={Precaución}]\small}{\end{tcolorbox}}

% ============================================================
% 8. CÓDIGO
% ------------------------------------------------------------
% listings Y NO minted: minted EXIGE --shell-escape Y Pygments,
% Y ESO ROMPE LA AUTONOMÍA DEL PROYECTO EN OTRA MÁQUINA.
% ============================================================
\lstset{
  basicstyle=\ttfamily\footnotesize,
  breaklines=true,
  columns=fullflexible,
  frame=none,
  showstringspaces=false,
  extendedchars=true,
  inputencoding=utf8
}

% ============================================================
% 9. COLOFÓN
% ------------------------------------------------------------
% LAS FUENTES Y LA VERSIÓN SON VALORES, NO TEXTO ESCRITO. SI
% UNA EDITORIAL REEMPLAZA LA ESTÉTICA Y CAMBIA LAS FUENTES,
% EL COLOFÓN SIGUE DICIENDO LA VERDAD SIN QUE NADIE SE
% ACUERDE DE EDITARLO.
% LA FECHA ES LA DE PUBLICACIÓN Y NO LA DEL DÍA: UNA MARCA DE
% TIEMPO HARÍA QUE CADA COMPILACIÓN PRODUJERA UN DIFF AUNQUE
% NADA HUBIERA CAMBIADO.
% ============================================================
% SI LA ESTÉTICA NO LAS DECLARA, EL COLOFÓN LO DICE EN VEZ DE DAR UN
% "Undefined control sequence" QUE NO EXPLICA NADA.
\providecommand{\gbfuentetexto}{(sin declarar)}
\providecommand{\gbfuentetitulos}{(sin declarar)}
\providecommand{\gbversion}{(sin version)}

% EL RECUADRO SE DEFINE COMO ESTILO Y NO EN LÍNEA: ASÍ UNA EDITORIAL
% QUE QUIERA OTRO ASPECTO REDEFINE gbcolofoncaja SIN TOCAR EL TEXTO,
% QUE ES LO QUE EL CONTRATO GARANTIZA.
\tcbset{gbcolofoncaja/.style={
  colback=black!7,
  colframe=black!50,
  boxrule=0.7pt,
  arc=3mm,
  left=8pt, right=8pt, top=8pt, bottom=8pt,
  boxsep=0pt}}

\newcommand{\gbTextoColofon}{%
  \begin{tcolorbox}[gbcolofoncaja]
    La composición tipográfica de este libro se realizó con
    gbpublisher versión~\gbversion, bajo un modelo de fuente única que
    garantiza la trazabilidad entre el original y cada una de sus
    salidas (HTML, ePub y PDF). Las familias tipográficas utilizadas son
    \gbfuentetexto\ para el texto e \gbfuentetitulos\ para los
    títulos. El proyecto se encuentra disponible en
    \url{https://github.com/albertomoyano/gbpublisher}.
  \end{tcolorbox}}

